% ACL format - Methods and Experimental Evaluation
% Overleaf version - no external .sty or .bst required

\documentclass[11pt]{article}
\usepackage[T1]{fontenc}
\usepackage{times}
\usepackage{latexsym}
\usepackage{amsmath}
\usepackage{booktabs}
\usepackage{graphicx}
\usepackage{multirow}
\usepackage{url}
\usepackage{caption}
\usepackage[a4paper, top=2.5cm, bottom=2.5cm, left=2.5cm, right=2.5cm, columnsep=0.6cm]{geometry}
\usepackage{multicol}

\usepackage{natbib}

\begin{document}
\begin{multicols}{2}

\section{Methods}
\subsection{Baseline System: CRF}
The baseline system is a CRF tagger that treats each PIO field independently. One model is trained for each of Participants (P), Interventions (I), and Outcomes (O), and the predictions are merged into a unified BIO tag sequence with label set $\{\text{B-P}, \text{I-P}, \text{B-I}, \text{I-I}, \text{B-O}, \text{I-O}, \text{O}\}$. When two fields simultaneously assign a non-O label to the same position, a fixed priority order of P $>$ I $>$ O resolves the conflict.

We build the feature set by combining several types of signals. First, we add orthographic features such as capitalization, all-uppercase, digit, and hyphen. Next, we extract prefix and suffix features that are two to four characters long. We also include the surface forms of the tokens before and after each word. Finally, we add domain keyword features that activate when a token matches a curated vocabulary for each field. For example, the vocabulary for P includes “patients”, for I it includes “placebo”, and for O it includes “mortality”.

\subsection{Axes of Experimentation}
The main design choice is between end-to-end extraction and a decomposed pipeline. In an end-to-end system, one model predicts all entity types together using the entire token sequence. In a decomposed system, a relevance filter first selects sentences, and span extraction is done only on those selected sentences.

The end-to-end approach assumes that joint training is beneficial since recognizing an intervention may provide context for locating the associated outcome. The decomposed approach assumes that filtering background sentences will reduce false positives and improve precision. Consequently, a trade-off exists: errors at the filtering stage propagate to the extraction stage and may lower recall.

\subsection{End-to-End BioBERT}
The end-to-end approach fine-tunes a single BioBERT model \cite{lee2020biobert} for multi-class token classification over the seven-label tag set. Sub-word tokenization is handled by aligning labels to the first sub-token. Continuation sub-tokens are assigned $-100$ and excluded from the loss.

We hypothesize that self-attention across the full input sequence will outperform the CRF, since the model can utilize broader sentence context when making each labeling decision. Training proceeds in two phases: firstly, the BERT encoder is frozen, and only the classification head is updated for four epochs. After that, all parameters are unfrozen, and the full model is fine-tuned for two additional epochs at a learning rate of $2 \times 10^{-5}$.

\subsection{Decomposed Pipeline}
First, we use a TF-IDF vectorizer with logistic regression to predict if a sentence contains a span for a specific field. We keep this approach simple because the relevance decision is a basic yes-or-no task, and a linear model works well for this purpose. A heavier model would only increase the computational cost without bringing any proportional benefit. In the second step, we fine‑tune three separate BioBERT models, one per field, and each model uses the output set $\{\text{O}, \text{B-X}, \text{I-X}\}$. We balance the training data by undersampling the negative sentences, so that the ratio of positive to negative sentences reaches 1:1. We hypothesize that relevance filtering will reduce false positives and, consequently, improve precision. Nevertheless, recall may decrease if the filter incorrectly discards relevant sentences.

\section{Experimental Evaluation}
\subsection{Performance Metrics}

All methods are evaluated using span-level precision, recall, and F1, calculated with the seqeval library \cite{nakayama2018seqeval}. A predicted span is a true positive only when its start index, end index, and entity type all exactly match the gold annotation \cite{nye2018corpus}. Results are reported per field and as the macro average, which treats each field equally.

Span-level F1 is strict about boundary errors. If a span has the correct entity type but is off by one token, it is counted as both a false positive and a false negative since there is no partial credit. This is especially important for Participants, where span boundaries in EBM-NLP are not always consistent.

\subsection{Experimental Setup}
The EBM-NLP dataset provides annotations for 4,457 training documents and 184 test documents.
The training set is split 90/10 for neural models. All experiments use a fixed random seed of 42 with no hyperparameter search. The Hugging Face Transformers library \cite{wolf2020transformers} is employed for all neural models. Batch sizes are 8 and 4 for the two training phases of the end-to-end model, and 8 for the per-field extractors. The maximum sequence length is 512 for the end-to-end model and 256 for the per-field extractors.

\subsection{Results}
Table \ref{tab:results} presents the span-level F1 and Table \ref{tab:pr} the precision and recall for each method. The end-to-end BioBERT achieves the highest macro F1 of 0.429, compared to 0.360 for the decomposed pipeline and 0.265 for the CRF.

% Tables must use \captionof inside multicols (table float not supported)
\begin{center}
\small
\begin{tabular}{lcccc}
\toprule
\textbf{Method} & \textbf{P} & \textbf{I} & \textbf{O} & \textbf{Macro} \\
\midrule
CRF        & 0.180 & 0.379 & 0.238 & 0.265 \\
End-to-End & 0.258 & 0.595 & 0.434 & 0.429 \\
Decomposed & 0.237 & 0.456 & 0.388 & 0.360 \\
\bottomrule
\end{tabular}
\captionof{table}{Span-level F1 on the EBM-NLP test set.}
\label{tab:results}
\end{center}

\begin{center}
\small
\begin{tabular}{llcc}
\toprule
\textbf{Method} & \textbf{Field} & \textbf{Precision} & \textbf{Recall} \\
\midrule
CRF        & P & 0.300 & 0.128 \\
           & I & 0.572 & 0.284 \\
           & O & 0.369 & 0.175 \\
\midrule
End-to-End & P & 0.253 & 0.262 \\
           & I & 0.608 & 0.583 \\
           & O & 0.456 & 0.414 \\
\midrule
Decomposed & P & 0.225 & 0.252 \\
           & I & 0.605 & 0.366 \\
           & O & 0.400 & 0.376 \\
\bottomrule
\end{tabular}
\captionof{table}{Span-level precision and recall per field.}
\label{tab:pr}
\end{center}

The CRF gives us a precision of 0.572 for Interventions, which is not bad at all, but then the recall figures come out as 0.128, 0.284, and 0.175 across the three fields. This is exactly what you would expect if keyword features are simply too rigid to capture clinical language as it is actually written. The end-to-end model substantially improves recall, particularly for Interventions (0.583) and Outcomes (0.414), with precision remaining competitive. The decomposed pipeline achieves precision close to the end-to-end model for Interventions (0.605 vs. 0.608), but recall drops to 0.366, giving an F1 of 0.456 compared to 0.595. Consequently, the precision improvement hypothesised in Section 3.4 is not clearly observed in the macro average.

Figure \ref{fig:comparison} shows the F1 by entity type and the precision-recall tradeoff across all three methods.
Figure \ref{fig:length} shows how F1 varies with span length for the CRF and decomposed pipeline, and confirms that longer spans are consistently harder to extract.

% Figures also use \captionof inside multicols
\begin{center}
\includegraphics[width=\columnwidth]{comparison.png}
\captionof{figure}{F1 by entity type and macro precision-recall tradeoff for all three methods.}
\label{fig:comparison}
\end{center}

\begin{center}
\includegraphics[width=\columnwidth]{f1_by_length.png}
\captionof{figure}{Span-level F1 by span length (tokens) for CRF and decomposed pipeline.}
\label{fig:length}
\end{center}

\subsection{Error Analysis}
To understand the failure modes, we examine false negative spans across both methods.
The CRF misses 579, 1055, and 1111 spans for P, I, and O, respectively (out of 664, 1473, and 1347 gold spans), whereas the decomposed pipeline reduces these to 497, 934, and 840, respectively.

Inspection of specific missed spans reveals two recurring patterns. Firstly, long multi-word spans are difficult for all methods: for example, the Interventions span Cholesteryl Ester Transfer Protein (CETP) inhibitor torcetrapib is missed by both, and F1 decreases as span length increases beyond five tokens (Figure \ref{fig:length}). Secondly, spans in sentence-initial position, such as Aquatic exercisers as an Interventions span, cause errors because these positions are not preceded by the context words the models rely on.

For Participants specifically, all methods perform poorly (with a maximum F1 score of 0.258). It is obvious from the error inspection that this is partially attributable to annotation inconsistency: EBM-NLP labels are crowd-sourced, and span boundaries are not always consistently defined, as seen in missed spans such as acute invasive diarrhea and Diabetic nephropathy patients. Consequently, annotation noise places an upper bound on the achievable F1 regardless of the method.

\subsection{Discussion}
The results show that contextual pre-training clearly improves performance compared to hand-crafted features. The end-to-end model scores 16.4 points higher than the CRF in macro F1, mainly due to better recall. This improvement happens because contextual embeddings can recognize surface forms that keyword features miss.

The decomposed pipeline does not achieve the expected precision improvement.
Although the relevance classifiers achieve F1 values of 0.770, 0.734, and 0.755 for the three fields, the errors they introduce compound with extraction errors, and the recall loss is not offset by a corresponding precision gain.
The reason is that the classifier occasionally filters out sentences containing spans in atypical language, and these spans are permanently missed before the extractor can recover them.
In future work, utilizing a lower classification threshold in the relevance stage may improve recall without substantially increasing computational cost.

\section*{References}
\bibliographystyle{acl_natbib}
\bibliography{references}

\end{multicols}
\end{document}